\section{Introducción}

El presente informe de implementación documenta el proceso de adopción, configuración y puesta en marcha del sistema de gestión en la nube basado en la plataforma \textbf{Odoo 19} en la empresa \textbf{Machu Picchu Exclusive Tours E.I.R.L.}, con el objetivo de centralizar la información crítica de clientes, profesionalizar la atención al turista extranjero y ordenar el flujo logístico--comercial asociado a la venta de paquetes turísticos de alta gama.

Machu Picchu Exclusive Tours E.I.R.L. es una \textit{empresa individual de responsabilidad limitada} dedicada principalmente a las \textbf{actividades de agencias de viajes} (CIIU 7911), con operaciones enfocadas en turismo receptivo de lujo. El 100\% de sus clientes son \textbf{turistas extranjeros}, y la gestión comercial se desarrolla de forma predominantemente virtual, a través de canales como \textit{WhatsApp}, recomendaciones directas y videollamadas personalizadas.

Previo al proyecto, la empresa no contaba con un sistema centralizado que integrara la información de los pasajeros (PAX), los itinerarios, los adelantos y saldos, y el seguimiento operativo de las reservas. La dispersión de datos en múltiples conversaciones de mensajería instantánea y archivos aislados dificultaba la trazabilidad, incrementaba el riesgo de errores en la emisión de boletos y consumía tiempo operativo valioso.

En este contexto, la implementación del sistema Odoo 19 se plantea como una solución estratégica para:

\begin{itemize}
    \item Centralizar la información de clientes y viajes en un CRM turístico seguro.
    \item Ordenar el flujo operativo de itinerarios mediante tableros Kanban.
    \item Controlar adelantos y saldos antes de realizar compras no reembolsables.
    \item Mejorar la eficiencia y la velocidad de respuesta al cliente internacional.
\end{itemize}

El presente documento describe, de manera estructurada, los antecedentes del proyecto, los objetivos perseguidos, el alcance funcional y organizacional, la metodología de trabajo, la situación inicial, las configuraciones realizadas en el sistema Odoo 19 y los resultados alcanzados tras la implementación.

